\chapter*{Notation and glossary}
\addcontentsline{toc}{chapter}{Notation and glossary}

Papers 1--7 write a permutation as a sequence of heights $h$, with the dot of column $i$ in row
$h_i$; Paper~8 writes it as $\pi$. All indices are zero-based.

\section*{Notation}
\begin{longtable}{@{}l >{\raggedright\arraybackslash}p{0.62\textwidth}@{}}
\toprule
Symbol & Meaning \\
\midrule
\endhead

$n$ & Order of a Costas array; the board is $n\times n$. Usually $n=32$. \\
$h$ & A permutation $h:\{0,\dots,n-1\}\to\{0,\dots,n-1\}$, written zero-based.
      $h_i$ is the \emph{height} (row) of the dot in column $i$. \\
$\Delta_{m,i}$ & $h_{i+m}-h_i$, the vertical difference at horizontal
      separation $m$ starting at column $i$. \\
$m$, $d$ & Horizontal separation (distance) between two columns. $d$ is used in
      solver contexts, $m$ in partition-counting contexts; they denote the same
      quantity. \\
$e$ & A signed vertical difference; $e\in\mathcal V$. \\
$g$ & An absolute vertical distance, $g=|e|$. \\
$\mathcal V$ & The signed palette $\{-(n-1),\dots,-1,1,\dots,n-1\}$,
      $|\mathcal V| = 2n-2$; for $n=32$, $|\mathcal V|=62$. \\
$L_m$ & $n-m$, the number of column pairs at separation $m$. \\
$\score(h)$ & Costas defect: the number of repeated $(d,\Delta)$ pairs,
      $\sum_{d,e}\max(0,\chi_{d,e}-1)$, over $1\le d\le n-1$ and nonzero $e$ with $|e|\le n-1$ (differences of a permutation are never zero). $\score(h)=0$ iff $h$ is Costas. \\
$\chi_{m,e}$ & $\#\{i: \Delta_{m,i}=e\}$, the occupancy of signed value $e$
      at separation $m$. \\
$y_{m,g}$ & $\chi_{m,g}+\chi_{m,-g}\in\{0,1,2\}$, magnitude occupancy. \\
$T_m$ & $\sum_i \Delta_{m,i}$, the telescoping row sum at separation $m$. \\
$N_q(r)$ & $\#\{e\in\mathcal V : e\equiv r \pmod q\}$, the signed residue
      capacity. \\
$\vtwo(x)$ & The 2-adic valuation of $x$; $\vtwo(0)=\infty$. \\
$D_i$ & Domain of the height variable $h_i$ during propagation. \\
$E_{i,j}$, $D_{ij}$ & Domain of the edge (difference) variable $c_{ij}=h_j-h_i$. \\
$k$ & Number of columns and rows deleted in an augmentation rung; the core has
      order $n-k$. \\
$K$ & A Costas core array of order $n-k$. \\
$C$, $R$ & The deleted column set and deleted row set, each of size $k$. \\
$\delta,\epsilon$ & Compressed displacement of a core pair: column distance and
      row difference \emph{within the core}. \\
$D_4$ & The dihedral group of order 8 acting on the square; the Costas
      symmetry group. \\
$W_1(p,g,c)$ & Welch construction, order $p-1$: $h_i = g^{i+c}\bmod p - 1$. \\
$W_2,W_3$ & $W_1$ with one or two corner dots deleted, orders $p-2$, $p-3$. \\
$G_2(q,\alpha,\beta)$ & Golomb construction over $\F_q$, order $q-2$:
      dot at $(i,j)$ iff $\alpha^i+\beta^j=1$. \\
$G_3,G_4,G_5$ & $G_2$ with corner dots deleted, orders $q-3,q-4,q-5$. \\
$C_{31}$ & The unique order-31 Costas array reachable by one-dot augmentation
      of the algebraic order-30 cores (one $D_4$ orbit). \\
$N_{32}$ & The known order-32 permutation of defect 2; the corner extension of
      $C_{31}$. \\
$\pi$ & In Paper~8, a permutation, written $\pi$ rather than $h$; transpose symmetry is
      $\pi=\pi^{-1}$. \\
$\Fix(\pi)$, $f$ & The fixed points of an involution $\pi$, and their number. \\
$F$, $R_F$ & A fixed-point set, and the region of Costas involutions whose fixed-point set is
      exactly $F$ (Paper~8). \\
$x_{ab}$, $h_{ij}$, $p_{s,x,\delta}$ & The edge, image and displacement-indicator variables of
      Paper~8's encoding of a region. \\
$Z(\pi)$ & In Paper~8, the number of large-large pairs: pairs $\{x,y\}$ with $|x-y|\ge16$ and
      $|\pi(x)-\pi(y)|\ge16$. \\
$k$ (Paper~8) & The number of split coordinates in a native split. In Papers~2 and~3, $k$ is the
      rung parameter instead. \\
$M(D,k)$ & In Paper~8, the number of messages of a native split with $k$ split coordinates and $D$
      free partners. \\
\bottomrule
\end{longtable}

\section*{Glossary}
\begin{description}[style=nextline,leftmargin=0pt,itemsep=2pt]
\item[Costas array] A permutation array of order $n$ in which all $\binom{n}{2}$ vectors connecting
pairs of dots are distinct. Equivalently, for every horizontal separation $m$, the vertical
differences $h_{i+m}-h_i$ are pairwise distinct.
\item[Displacement] The vector $(j-i,\,h_j-h_i)$ between the dots of columns $i<j$.
\item[Defect, score] The number of repeated displacement vectors, $\score(h)$. A Costas array has
defect 0; the known order-32 near-solution $N_{32}$ has defect 2.
\item[Dihedral orbit, equivalence class] The set of arrays obtained from one array under the eight
symmetries of the square, $D_4$. Censuses count arrays both individually and by class.
\item[Census] A complete list of the Costas arrays of one order, obtained by exhaustive enumeration.
Complete censuses exist through order 29.
\item[Census-complete, catalogue-relative] A result is census-complete if it quantifies over a
complete census, and catalogue-relative if it quantifies only over a generated catalogue, such as
the 664 order-30 arrays of Paper~2, at an order that has not been enumerated.
\item[Welch and Golomb constructions] The two principal algebraic constructions of Costas arrays,
over prime fields and finite fields respectively; see the notation table.
\item[Circular and singly periodic sequences] A Costas sequence $a_0,\dots,a_{n-1}$ with values in
$\{1,\dots,n\}$ is \emph{circular} if, for every lag $k$, the circular differences
$a_{(i+k)\bmod n}-a_i$ are distinct modulo $n+1$: its dots, read in $\Z_n\times\Z_{n+1}$, have
distinct displacements. It is \emph{singly periodic} if every cyclic shift of its terms is again
Costas. Circularity lies between single periodicity and Moreno's shifting property. Circular
sequences exist only when $n+1$ is prime (Golomb and Moreno, 1996) and are exactly the Welch
sequences (Muratovi\'c-Ribi\'c, Pott, Thomson and Wang, 2015); every circular sequence is an
ordinary Costas array, but not conversely. Arrays required to stay Costas when extended periodically
in both directions are a further notion, studied separately. See Paper~1.
\item[Corner deletion, corner addition] Removing or adding a dot at a corner of the board, together
with its row and column. The compression is then a translation, so these are the deletions that
preserve the Costas property in general.
\item[Sporadic array] A Costas array not accounted for by a recognised generated or emergent
construction. The category shrinks whenever a new construction is found. Order 28 is the first
order above 5 with no sporadic arrays, and order 29 has none either.
\item[Emergent array] An array produced by a semi-empirical construction rather than an algebraic
family, such as the 16 order-29 arrays produced by the Rickard--Welch operation.
\item[Augmentation rung] The exhaustive search over all permutations of order $n$ formed by placing
a Costas core of order $n-k$ in a compressed sub-board and adding $k$ dots; see the note on software
engineering.
\item[Deletion minor] The permutation array left after deleting $k$ columns, $k$ rows and the $k$
dots they share from a larger permutation array, and closing up the gaps: the inverse of a rung-$k$
augmentation.
\item[Core-first pruning] In a rung search, testing the placed core's own displacements before
considering the added dots. Exact for defect 0, since a repeated displacement among a subset of the
dots is a repeated displacement of the whole.
\item[Involution, fixed point, transposition] A permutation equal to its inverse; a coordinate $c$
with $\pi(c)=c$, whose dot lies on the main diagonal; a pair $(a\,b)$ with $\pi(a)=b$ and $\pi(b)=a$.
\item[Geometric and transpose symmetry] An array is geometrically symmetric if a non-identity motion
of the square fixes it, and transpose-symmetric if the main-diagonal reflection does. For orders of
at least 3 the two notions agree up to the choice of diagonal (Paper~8). The literature calls
transpose-symmetric arrays \emph{diagonal}.
\item[Anti-reflective and consecutive arrays] Symmetry types studied in the literature that are not
motions of the square. They are mentioned for comparison only.
\item[Levels of assurance] A \emph{validated} computation agrees with independent implementations,
regression tests or known answers; an \emph{audited} computation has its coverage and execution
records independently checked; a \emph{machine-checked witness} is an object, such as a Costas
permutation or a Lean theorem, checked independently; a \emph{proof-certified} unsatisfiable verdict
has a checked DRAT or LRAT proof; an \emph{uncertified} one rests on the solver's exhaustion.
\item[Region] In Paper~8, the set of Costas involutions with a given fixed-point set.
\item[Native split, message, unsplit branch] Paper~8's distributed decomposition fixes the partners
of $k$ chosen coordinates; each assignment of partners is a message, and a message in which two
chosen coordinates are paired with each other is an unsplit branch.
\item[Seed-refuted] In Paper~8, a pairing that repeats a displacement with the fixed dots alone.
\item[Case] In Paper~8, a region with one transposition fixed.
\item[Gate] An acceptance criterion frozen before the run it judges; see the note on software
engineering.
\item[CDCL, DRAT and LRAT] Conflict-driven clause learning, the algorithm of modern SAT solvers; and
machine-checkable proof formats for unsatisfiability. Without such a proof, an unsatisfiable verdict
is a computational observation rather than a certificate.
\item[Verdicts] Paper~8 records \st{SAT-VALIDATED} (a model that passed independent checking),
\st{UNSAT-UNCERTIFIED} (normal solver exhaustion, no checked proof), \st{TIMEOUT} and \st{ERROR}. A
region closed by an \st{UNSAT-UNCERTIFIED} verdict is \emph{search-closed}.
\end{description}
